\chapter{Временная граница барионного маршрута и возврат к общему полевому мосту}
\label{ch:v8-temporary-boundary-and-retrospective-return}

Последовательность барионных проверок не привела к локальному
шеститочечному родителю. Это не отменяет построенный связный оператор
$W_3$ и не означает провала всей программы. Получен более узкий вывод:
при существующем действии локальная беспроизводная ветвь достигла
временного тупика.

\section{Почему локальное продолжение остановлено}

Три последние возможности исчерпаны независимо:
\begin{equation}
 \begin{aligned}
 M=\frac m2I_{21}&:\quad \Gamma M+M\Gamma=m\Gamma,\\
 D_{\text{неч}}&:\quad
 \operatorname{supp}c^D\subset TTT\oplus TGG,\\
 \mathbb F_{\text{пространство-время}}&:\quad
 S^{(3)}=0\quad\text{при }dA=d\Phi=0.
 \end{aligned}
 \label{eq:v8-temporary-boundary-three-local-obstructions}
\end{equation}
Центральный фон имеет неверную чётность, нечётный фон поддержан в каналах,
дополнительных к $d_{abc}$, а допустимые пространственно-временные
$TTG/GGG$-вершины содержат производные и исчезают на постоянном
нулеимпульсном срезе. Поэтому следующий локальный полином отличался бы от
уже проверенных вариантов только новым предположением или свободным
коэффициентом.

Настоящее продолжение барионной линии потребовало бы нового объекта:
\begin{equation}
 K_{(3)}(p_1,p_2,p_3;p'_1,p'_2,p'_3),
 \label{eq:v8-temporary-boundary-nonlocal-six-point-kernel}
\end{equation}
то есть производного или нелокального шеститочечного ядра с явной
импульсной зависимостью. Такое ядро пока не следует из принятого
пространства полей и потому не вводится только ради продолжения вычислений.

\section{Что накоплено к моменту остановки}

Одновременно Том VIII построил два полных объекта, которых не существовало
при постановке ранних задач:
\begin{equation}
 \mathfrak F_{\text{поле}}
 =(\mathbb C^{15})_{\mathbb R}\oplus\mathbb R^{12},
 \qquad
 \mathcal N_{\text{шум}}
 =\operatorname{span}_{\mathbb R}\{F_a\}_{a=1}^{42}.
 \label{eq:v8-temporary-boundary-two-complete-carriers}
\end{equation}
Оба пространства имеют вещественную размерность $42$. Для шумового репера
уже известны невырожденная метрика следа $K_{ab}=\Tr(F_aF_b)$ и её точная
обратная. Раньше совпадение размерностей было лишь условием допустимости;
теперь можно построить и проверить само отображение.

\section{Причина возврата к старому вопросу}

Поэтому дальнейшая линия меняется не из-за произвольного выбора темы, а
из-за изменения доказательных предпосылок. Естественный кандидат равен
\begin{equation}
 \mathcal J(\delta A,\delta B_s,\delta B_t)
 =\begin{pmatrix}
   \delta B_s&\delta A^*\\
   \delta A&\delta B_t
  \end{pmatrix}.
 \label{eq:v8-temporary-boundary-field-noise-map}
\end{equation}
Нужно проверить его полный ранг, калибровочное сплетение и обратный перенос
метрики
\begin{equation}
 G_{\text{поле}}=\mathcal J^*K\mathcal J,
 \qquad
 R_{\text{поле}}=\mathcal J^{-1}K^{-1}(\mathcal J^{-1})^*.
 \label{eq:v8-temporary-boundary-pulled-metrics}
\end{equation}
Именно здесь старый вопрос о связи родительского гессиана с мобильностью
полного квантово-марковского процесса впервые стал вычислимым на полном
носителе.

\begin{remark}
Возврат не переоткрывает закрытые численные формулы, абсолютную единицу
времени или электромагнитную барионную величину. Он использует накопленную
структуру для проверки более раннего фундаментального разрыва между полем
и процессом.
\end{remark}

Следующая содержательная глава должна поэтому исследовать отображение
$\mathcal J$ и обратную следовую метрику. Если они не согласуются с полевым
гессианом, кинематический мост всё равно будет построен, а необходимость
дополнительного закона флуктуации--диссипации станет строгим результатом.